% ══════════════════════════════════════════════════════════════════════════════
%  EXPERIMENT 3 — FAST AND SLOW EVAPORATION                      (Lesson 9, Term 2 Week 3)
% ══════════════════════════════════════════════════════════════════════════════
\experimentdiv{Fast and Slow Evaporation}{Evaporating to dryness safely \textbullet{} Observing crystals with a hand lens \textbullet{} Linking cause and effect}{Lesson 9 \textbullet\ Term 2, Week 3 \textbullet\ Lab}\label{sec:3}
\lessoncodes{VC2S8U05, U07 \textbullet{} VC2S8I02, I05, I07, I08}

\vocabhead
\begin{ctab}
\begin{tabular}{|SW{4.2cm}|W{11.6cm}|}
  \hline\hd & \thead{Word} & \thead{Definition (copy from the slide)} \\ \hline
  & \scaf{\textbf{Evaporation}} & \scaf{A liquid changing to a gas at its surface, below its boiling point.} \\ \hline
  & \textbf{Crystal} & \\ \hline
  & \textbf{Crystallisation} & \\ \hline
  & \textbf{Residue} & \\ \hline
  & \textbf{Cause and effect} & \\ \hline
\end{tabular}
\end{ctab}

\begin{lipanel}
\inquiryli{Heat a solution safely and draw what you observe with a hand lens.}
\understandli{Explain how the rate of evaporation affects the size of the crystals that form.}
\sctitle
\begin{itemize}[leftmargin=1.4em, itemsep=0pt, topsep=1pt]
  \item I can heat a solution safely and stop at the right time.
  \item I can draw a labelled diagram of the set-up and of the crystals.
  \item I can explain why slow evaporation makes larger crystals.
\end{itemize}
\end{lipanel}

\begin{skillsfocus}
\focus{PC3}{Heat an evaporating basin safely.}
\focus{PC6}{Observe the crystals with a hand lens.}
\focus{CO2}{Draw the set-up and the crystals.}
\focus{PA6}{Link the rate of evaporation to crystal size.}
\focus{EV5}{Suggest a question for a further investigation.}
\end{skillsfocus}

% ─── TECHNIQUE ────────────────────────────────────────────────────────────────
\exppart{Technique spotlight}{Blue}

\begin{technique}{Evaporating a solution with a Bunsen burner}{PC3}
\begin{enumerate}[leftmargin=1.6em, itemsep=1pt, topsep=3pt]
  \item Stand the evaporating basin on a gauze mat on a tripod, with a heatproof mat under the Bunsen.
  \item Use a \textbf{blue} flame, about half open. Stand up, wear safety glasses, and tie back long hair.
  \item Heat gently. When the first crystals appear at the edge, \textbf{turn the Bunsen off}. The hot basin finishes the job; heating to dryness makes the salt spit out.
  \item Leave the basin to cool on the gauze. Move it only with tongs once it has cooled.
\end{enumerate}
\end{technique}

\begin{workedex}{Drawing a scientific diagram}
A scientific diagram is a 2D side view drawn in pencil with a ruler. Every label sits to the side of the drawing, joined to its part by a straight line that does not cross another line.

\smallskip
\begin{center}
\begin{tikzpicture}[scale=0.62]
  \begin{scope}[shift={(0,0)}]\eqBunsen\end{scope}
  \begin{scope}[shift={(-0.4,2.4)}]\eqTripod\end{scope}
  \draw[eq] (-0.6,4.52) -- (2.4,4.52);
  \begin{scope}[shift={(-0.4,4.52)}]\eqBasin\end{scope}
  \fill[black!15] (-0.3,4.9) .. controls (0.2,4.6) and (1.6,4.6) .. (2.1,4.9) -- cycle;
  \eqlab{(2.0,5.3)}{(3.6,5.8)}{evaporating basin}
  \eqlab{(2.3,4.52)}{(3.6,4.7)}{gauze mat}
  \eqlab{(2.0,3.5)}{(3.6,3.5)}{tripod}
  \eqlab{(1.0,1.6)}{(3.6,1.8)}{Bunsen burner}
\end{tikzpicture}
\end{center}
\end{workedex}

\begin{wedo}{Spot the mistakes}
A student drew the set-up in pen, shaded the flame, and wrote the labels on the drawing. List what they should change.
\answerbox{3.5cm}
\end{wedo}

% ─── PLAN AND SAFETY / METHOD ────────────────────────────────────────────────
\exppart{Plan, safety and method}{Teal}

\textbf{Aim:} To compare the crystals that form when a salt solution evaporates quickly and slowly.

\task Predict which method will make the larger crystals. Give a reason. \skr{QP1}
\answerbox{4cm}

\subsubsection*{Equipment}
Saturated salt solution \textbullet{} 10\,mL measuring cylinder \textbullet{} evaporating basin \textbullet{} watch glass \textbullet{} Bunsen burner, tripod, gauze mat and heatproof mat \textbullet{} tongs \textbullet{} hand lens \textbullet{} safety glasses

\begin{infobox}{Safety}
Hot glass and hot basins look exactly like cold ones: touch them only with tongs. Turn the Bunsen off as soon as crystals appear, so that hot salt does not spit. Keep the gas off when you are not using it.
\end{infobox}

\subsubsection*{Method}
\textbf{Fast evaporation}
\begin{enumerate}[itemsep=2pt]
  \item Pour 10\,mL of salt solution into the evaporating basin and set up the Bunsen as in the technique box.
  \item Heat gently until crystals start to form at the edge. Turn off the gas and leave the basin to cool.
\end{enumerate}
\textbf{Slow evaporation}
\begin{enumerate}[itemsep=2pt, resume]
  \item Pour 5\,mL of salt solution onto a watch glass and label it with your names.
  \item Leave it on the windowsill until next lesson.
\end{enumerate}
\textbf{Both}
\begin{enumerate}[itemsep=2pt, resume]
  \item Look at each set of crystals through the hand lens. Record your observations and draw them.
\end{enumerate}

% ─── RESULTS ─────────────────────────────────────────────────────────────────
\exppart{Results}{Amber}

\task Draw a scientific diagram of the fast-evaporation set-up. \skr{CO2}
\drawbox{7cm}{}

\subsubsection*{Table 3.1: Crystals from fast and slow evaporation \skr{PC6}}
\begin{ctab}
\begin{tabular}{|SW{3.4cm}|W{3.4cm}|W{3.6cm}|W{4.6cm}|}
  \hline\hd & \thead{Method} & \thead{Time taken} & \thead{Size of crystals} & \thead{Shape and appearance} \\ \hline
  & \scaf{Example: sugar, fast} & \scaf{5 minutes} & \scaf{Too small to see one by one} & \scaf{White crust, like powder} \\ \hline
  & Fast (Bunsen) & & & \\ \hline
  & Slow (windowsill) & & & \\ \hline
\end{tabular}
\end{ctab}

\needspace{14\baselineskip}
\task Draw what you saw through the hand lens. \skr{CO2}

\noindent
\begin{minipage}[t]{0.48\linewidth}\drawbox{6cm}{}\nameline{Fast evaporation}\end{minipage}\hfill
\begin{minipage}[t]{0.48\linewidth}\drawbox{6cm}{}\nameline{Slow evaporation}\end{minipage}

% ─── ANALYSIS ────────────────────────────────────────────────────────────────
\exppart{Analysis and evaluation}{Coral}

\task Was your prediction correct? Use your observations to answer.
\answerbox{3.5cm}

\needspace{12\baselineskip}
\task Complete the sentence: \emph{The faster a solution evaporates, the \dots\ the crystals.} Then use the particle model to explain why. \skr{PA6}
\answerbox{5cm}

\needspace{10\baselineskip}
\task The salt crystals were white, but there was no sign of water in the basin at the end. Where did the water go? Name the change of state.
\answerbox{3.5cm}

\needspace{10\baselineskip}
\task Salt farms in South Australia make salt by letting seawater evaporate in large shallow ponds. Explain why the ponds are wide and shallow rather than deep.
\answerbox{4cm}

\needspace{10\baselineskip}
\task Write one question about crystals that you could test in another experiment. \skr{EV5}
\answerbox{3cm}

\begin{yr8box}
\begin{keepext}
\exttask A student says, ``The crystals were bigger on the windowsill because the windowsill was sunny.'' Is this a cause, or only something that happened at the same time? Plan how you could test it.
\answerbox{5cm}
\end{keepext}
\end{yr8box}

% ─── SKILL CHECK ─────────────────────────────────────────────────────────────
\exppart{Skill check}{NavyBlue}

\begin{skillcheck}
\scrow{PC3}{I heated the solution safely and stopped at the right time.}
\scrow{PC6}{I observed the crystals closely with a hand lens.}
\scrow{CO2}{I drew a scientific diagram of the set-up and the crystals.}
\scrow{PA6}{I can link the rate of evaporation to the size of the crystals.}
\scrow{EV5}{I can suggest a question for a further investigation.}
\end{skillcheck}

\task Which skill will you focus on in Experiment 4? Explain what you will do differently.
\answerbox{5cm}
